\documentclass{scrartcl}
\usepackage[latin1]{inputenc}
\usepackage[T1]{fontenc}
\usepackage[ngerman]{babel}
\usepackage{amsmath}
\usepackage{amssymb}
\usepackage{icomma}
%\usepackage[dvips]{graphicx}
%\usepackage{floatflt}
%\usepackage{enumitem}
%\usepackage{babel}
\usepackage{blindtext}
%\usepackage{showframe}
\usepackage{calc}
\usepackage{wrapfig}
\def\BILD{\rule{0.4\textwidth}{4cm}}

\usepackage{graphicx}
\usepackage{placeins}
\usepackage{multirow}
\usepackage{subfig}
\usepackage{url}

\renewcommand{\topfraction}{.85}
\renewcommand{\bottomfraction}{.7}
\renewcommand{\textfraction}{.15}
\renewcommand{\floatpagefraction}{.66}
\renewcommand{\dbltopfraction}{.66}
\renewcommand{\dblfloatpagefraction}{.66}
\setcounter{topnumber}{9}
\setcounter{bottomnumber}{9}
\setcounter{totalnumber}{20}
\setcounter{dbltopnumber}{9}
\setlength{\intextsep}{0cm plus1cm minus1cm}
\pdfminorversion = 5
\usepackage{setspace}
\onehalfspacing

\begin{document}
\title{Versuch 4: Bestimmung der spezifischen Ladung $\frac{e}{m}$ des
Elektrons}

\date{15.11.2010}

\author{Gruppe 5a: Gia-Danh Lam, Nils Haldenwang}

\maketitle
\tableofcontents

\newpage

\section{Einleitung}
\label{sec:einleitung}

Als \emph{spezifische Elektronenladung} bezeichnet man das Verh�ltnis
$\frac{e}{m}$ der Ladung \emph{e} eines Elektrons zu dessen Masse
\emph{m} . Die Konstante spielt insbesondere in der Atomphysik und den
angrenzenden Gebieten eine wichtige Rolle. Aus der \emph{spezifischen
  Elektronenladung} l�sst sich mit Kenntnis der \emph{Elementarladung
  e} die Masse eines Elektrons bestimmen.

Es existieren zahlreiche Verfahren zur Messung der \emph{spezifischen
  Elektronenladung}. Im folgenden werden zwei eher einfache
Messmethoden zur Anwendung gebracht. Zun�chst bestimmen wir die
Konstante aus der Kreisbahn von Elektronen im homogenen Magnetfeld. Im
zweiten Versuch wird ein Kr�ftegleichgewicht zwischen
Elektrostatischer Kraft und Lorentzkraft hergestellt und daraus der
gesuchte Wert berechnet.

Die Elektronenstrahlen werden in beiden Versuchen mittels einer
Gl�hkathode, einer Anode und eines Wehneltzylinders erzeugt. Der
Wehneltzylinder wurde 1902 / 1903 von Arthur Wehnelt entwickelt. Er
ist in direkter N�he der Gl�hkathode angebracht und ihr gegen�ber auf
einem negativen elektrischen Potential. Durch Einstellen dieser
Spannung l�sst sich die Intensit�t des Elektronenstrahls kalibrieren,
da langsame Elektronen das Potential nicht mehr �berwinden k�nnen.
Weiterhin werden Elektronen auf Bahnen fern der Strahlachse durch das
negative Potential der Zylinderwand abgesto�en und so der
Elektronenstrahl st�rker geb�ndelt.

Der Leuchteffekt des Strahles kommt zu Stande durch St��e der
Elektronen mit dem umgebenden Gas im Fadenstrahlrohr. Die Gasteilchen
werden durch den Sto� in einen angeregten Zustand versetzt und
strahlen Teile dieser Energie als sichtbares Licht wieder ab.
Insbesondere leuchten hier \emph{nicht} die Elektronen selbst!

Nach Einschalten der Heizspannung $U_{h}$ und Durchlaufen der
Beschleunigungsspannung $U_{a}$ haben die Elektronen eine kinetische
Energie von $E_{kin} = e \cdot U_{a}$. Tats�chlich sollte man hier
aber auch noch die thermische Geschwindigkeit ber�cksichtigen, welche die
Elektronen beim Austritt aus dem Kathodenmaterial erhalten, sodass die
effektive Beschleunigungsspannung
\begin{equation}
  \label{eq:1}
  U_{a effektiv} = U_{a} + \frac{k_B T}{2 e}
\end{equation}
verwendet werden sollte. Dabei ist $k_B$ die Boltzmannkonstante und
$T$ die Temperatur.


\subsection{Elektronen im homogenen Magnetfeld}
\label{sec:elektr-im-homog}
Die Magnetfelder werden bei beiden Messungen mittels eines
Helmholtzspulenpaares erzeugt. Diese Spulenanordnung wurde vom
deutschen Physiker Hermann von Helmholtz entwickelt. Sie besteht aus
zwei parallel auf gleicher Achse platzierten, d�nnen Spulen mit dem
Radius $R$. Durch die �berlagerung der Felder beider Spulen ergibt
sich in der N�he der Achse ein nahezu homogenes Feld, welches sich
recht gut analytisch berechnen l�sst:

\begin{equation}
B= K\cdot I_S= \mu_0\cdot \left(\frac{4}{5}\right)^{\frac{2}{3}}\frac{n\cdot I_S}{R}
\label{eq: Magnetfeld}
\end{equation}

Dabei beschreibt $n$ die Anzahl der Windungen, $R$ den Radius der
Spulen und $I_S$ den Spulenstrom.

Tats�chlich ist aber auch in dieser Anordnung das Magnetfeld nicht
absolut homogen. Die Inhomogenit�ten h�ngen direkt vom Abstand $r$ der
untersuchten Kreisbahn zur Achse der Spulen ab. Sie lassen sich mit
folgender Formel in Abh�ngigkeit von $r$ berechnen:

\begin{equation}
B(r)= \frac{\mu_0 \cdot n\cdot I_S}{2}R^2\left[\left(R^2+\left(r+\frac{R}{2}\right)^2\right)^{-\frac{3}{2}}+\left(R^2+\left(r-\frac{R}{2}\right)^2\right)^{-\frac{3}{2}}\right]
\label{eq: Inhomogen}
\end{equation}

Der fokussierte Elektronenstrahl wird nun durch dieses Magnetfeld
geleitet. Aufgrund der senkrecht zur Flugrichtung der Elektronen
wirkenden Lorentzkraft $F_L = e \cdot v \cdot B$ werden die Elektronen
auf eine Kurvenbahn abgelenkt. Sofern der Elektronenstrahl nicht aus
dem Magnetfeld hinausl�uft entsteht ein Kr�ftegleichgewicht der
Loretznkraft $F_L$ und der Zentripetalkraft $F_Z = \frac{m \cdot
  v^2}{r}$. Dies f�hrt dazu, dass sich die Elektronen auf einer
Kreisbahn bewegen. Mit Kenntnis der Beschleunigungsspannung $U_a$
l�sst sich durch Anwendung des Energieerhaltungssatzes $\frac{m}{2}
v^2 = e \cdot U_a$ die Geschwindigkeit aus dem Kr�ftegleichgewicht
eliminieren und man erh�lt die folgende Gleichung zur Bestimmung von
$\frac{e}{m}$:

\begin{equation}
  \label{eq:2}
  \frac{e}{m} = \frac{2 \cdot U_a}{B^2 \cdot r^2}
\end{equation}

Dabei bezeichnet $r$ den Radius der Kreisbahn. Das Magnetfeld $B$
berechnet sich nach Gleichung \eqref{eq: Magnetfeld}.


\subsection{Elektronen im gekreuzten homogenen elektrischen und magnetischen Feld}
\label{sec:elektr-im-gekr}

Sendet man einen wie in Abschnitt \ref{sec:einleitung} beschriebenen
Elektronenstrahl durch das homogene elektrische Feld $E_k$ eines
Plattenkondensators, so werden die Elektronen aufgrund der
elektrostatischen Kraft des Feldes auf eine parabelf�rmige Bahn in
Richtung der positiven Kondensatorplatte gezwungen. Diese Bahnkurve
kann durch folgende Gleichung beschrieben werden:

\begin{equation}
  \label{eq:3}
  y = \frac{1}{2} \frac{e}{m} \frac{E_k}{v_x^2} \cdot x^2
\end{equation}

Wird nun ein wie Abschnitt \ref{sec:elektr-im-homog} beschriebenes
Magnetfeld $B$ senkrecht zum elektrischen Feld $E_k$ hinzugenommen und
ein Kr�ftegleichgewicht zwischen der elektrostatischen Kraft und der
Lorentzkraft hergestellt, so kann man daraus die Geschwindigkeit $v_x$ der
Elektronen in x-Richtung bestimmen:

\begin{equation}
  \label{eq:4}
  v_x = \frac{E_k}{B}
\end{equation}

Man bezeichnet diese Anordnung als Wienfilter. Wilhelm Wien
entwickelte sie einst, um einen Elektronenstrahl aus Elektronen mit
einer bestimmten Geschwindigkeit zu erzeugen. Nur jene Elektronen mit
der Geschwindigkeit $v = v_x$ bleiben auf der waagerechten Bahn, alle
anderen werden abgelenkt und lassen sich mit einer Lochblende leicht
abfangen.
F�r unseren Versuch ist allerdings nur die Bestimmungsgleichung f�r
$\frac{e}{m}$ interessant, welche man durch einsetzen von Gleichung
\eqref{eq:4} in Gleichung \eqref{eq:3} erh�lt:

\begin{equation}
  \label{eq:5}
  y = \frac{1}{2} \frac{e}{m} \frac{B^2}{E_k} \cdot x^2
\end{equation}

Das Magnetfeld $B$ wird wieder nach Gleichung \eqref{eq: Magnetfeld}
bzw. Gleichung \eqref{eq: Inhomogen} berechnet.

\section{Bestimmung von $\frac{e}{m}$ aus der Kreisbahn im Magnetfeld}
\label{sec:best-von-frac}

\subsection{Versuchsaufbau}
\label{sec:versuchsaufbau}

\begin{figure}[htb!]
  \centering
  \includegraphics[width=12cm]{pics/dummypic}
  \caption{Schaltskizze des Versuchsaufbaus.}
  \label{fig:Schaltskizze-1}
\end{figure}

\subsection{Durchf�hrung}
\label{sec:durchfuhrung}

\subsection{Systematische Fehlerquellen}
\label{sec:fehlerquellen}

\subsection{Ergebnis}
\label{sec:ergebnis}

\section{Messungen mit elektrischem und magnetischem Feld}
\label{sec:mess-mit-elektr}

\subsection{Versuchsaufbau}
\label{sec:versuchsaufbau-1}

\subsection{Durchf�hrung}
\label{sec:durchfuhrung-1}

\subsection{Ergebnis}
\label{sec:ergebnis-1}

\end{document}